\subsection{Extension of a valuation and completion}

DEFINITION 3. Let K be a field, v a valuation on K and L an extension of K. A family \((v_{t}^{\prime})_{t\in I}\) of valuations on L which extend v and such that every valuation on L extending v is equivalent to a unique \(v_{t}^{\prime}\) is called a complete system of extensions of v to L.

PROPOSITION 2. Let K be a field, v a valuation on K, \(\hat{K}\) the completion of K with respect to v, \(\hat{v}\) the continuous extension of v to \(\hat{K}\) and L a finite extension of K of degree n.

\begin{enumerate}
\def\labelenumi{(\alph{enumi})}
\tightlist
\item
  Let \(v'\) be a valuation on \(\mathbf{L}\) extending \(v\) ; let \(\hat{\mathbf{L}}_{v'}\) denote the completion of \(\mathbf{L}\) with respect to \(v'\) and \(\hat{v}'\) the continuous extension of \(v'\) to \(\hat{\mathbf{L}}_{v'}\) ; identifying \(\hat{\mathbf{K}}\) with the closure of \(\mathbf{K}\) in \(\hat{\mathbf{L}}_{v'}\) ,
\end{enumerate}

\begin{enumerate}
\def\labelenumi{(\arabic{enumi})}
\setcounter{enumi}{3}
\tightlist
\item
\end{enumerate}

\[
e (\hat {v} ^ {\prime} / \hat {v}) = e (v ^ {\prime} / v), \quad f (\hat {v} ^ {\prime} / \hat {v}) = f (v ^ {\prime} / v),\tag{5}
\]

\[
[ \hat {\mathbf {L}} _ {v ^ {\prime}}: \hat {\mathbf {K}} ] \leqslant n,\tag{6}
\]

\[
e (v ^ {\prime} / v) f (v ^ {\prime} / v) \leqslant [ \hat {\mathbf {L}} _ {v ^ {\prime}}: \hat {\mathbf {K}} ].
\]

\begin{enumerate}
\def\labelenumi{(\alph{enumi})}
\setcounter{enumi}{1}
\tightlist
\item
  Every set of pairwise independent valuations on \(\mathbf{L}\) extending a non-improper valuation \(v\) is finite. Let \(v_{1}',\ldots ,v_{s}'\) denote pairwise independent valuations on \(\mathbf{L}\) extending v such that every valuation on L extending v is dependent on one of the \(v_{i}^{\prime}\) ; let \(L_{i}\) be the field L with the topology defined by \(v_{i}^{\prime}\) and \(\hat{L}_{i}\) its completion; we write \(n_{i} = [\hat{L}_{i}: \hat{K}]\) . Then the canonical mapping
\end{enumerate}

\[
\phi : \hat {\mathbf {K}} \otimes_ {\mathbf {K}} \mathbf {L} \rightarrow \prod_ {i = 1} ^ {s} \hat {\mathbf {L}} _ {i}
\]

(extending by continuity the diagonal mapping \(\mathbf{L} \to \prod_{i=1}^{s} \mathbf{L}_i\) ) is surjective, its kernel is the Jacobson radical of \(\hat{\mathbf{K}} \otimes_{\mathbf{K}} \mathbf{L}\) and

\[
\sum_ {i = 1} ^ {s} n _ {i} \leqslant n.\tag{7}
\]

Let us first prove (a). Suppose that \(v\) is not improper. As \(v\) and \(\hat{v}\) (resp. \(v'\) and \(\hat{v}'\) ) have the same order group and the same residue field (§ 5, no. 3, Proposition 5 (b) and (f)), (4) holds. We deduce (6) from it by means of Lemma 2. Finally the vector sub- \(\hat{K}\) -space of \(\hat{L}_{v'}\) generated by \(L\) is closed (§ 5, no. 2, Corollary to Proposition 4) and everywhere dense and hence equal to \(\hat{L}_{v'}\) ; this shows (5).

We now pass to (b). We may still assume that \(v\) is not improper. Let \((v_1', \ldots, v_r')\) be any finite family of pairwise independent valuations on \(L\) extending \(v\) . The image of \(L\) in \(\prod_{i=1}^{r} L_i\) under the diagonal mapping is everywhere dense (§ 7, no. 2, Theorem 1) and \(\prod_{i=1}^{r} L_i\) is dense in \(\prod_{i=1}^{r} \hat{L}_i\) . Hence the canonical image of \(\hat{K} \otimes_K L\) in \(\prod_{i=1}^{r} \hat{L}_i\) is everywhere dense. On the other hand this image is a vector sub- \(\hat{K}\) -space of \(\prod_{i=1}^{r} \hat{L}_i\) ; as \(\prod_{i=1}^{r} \hat{L}_i\) is of finite dimension over \(\hat{K}\) by (5), the image of \(\hat{K} \otimes_K L\) is closed (§ 5, no. 2, Corollary to Proposition 4) and hence equal to \(\prod_{i=1}^{r} \hat{L}_i\) . As the dimension of \(\hat{K} \otimes_K L\) over \(\hat{K}\) is \(n, \sum_{i=1}^{r} n_i \leqslant n\) . This shows in particular that the integer \(r\) is bounded above by \(n\) and shows the first assertion of (b).

We now take \((v_{1}',\ldots ,v_{s}^{\prime})\) as in the statement. The fact that

\[
\phi : \hat {\mathbf {K}} \otimes_ {\mathbf {K}} \mathbf {L} \rightarrow \prod_ {i = 1} ^ {s} \hat {\mathbf {L}} _ {i}
\]

is surjective and relation (7) have already been shown. It remains to verify that the kernel \(\mathfrak{n}\) of \(\phi\) is the Jacobson radical \(\mathfrak{r}\) of \(\hat{\mathbf{K}}\otimes_{\mathbb{K}}\mathbf{L}\) . As \(\prod_{i = 1}\hat{\mathbf{L}}_i\) is semi-simple, \(\mathfrak{r}\subset \mathfrak{n}\) . On the other hand, for every maximal ideal \(\mathfrak{m}\) of \(\hat{\mathbf{K}}\otimes_{\mathbb{K}}\mathbf{L}\) , the quotient field \(\mathbf{L}(\mathfrak{m}) = (\hat{\mathbf{K}}\otimes_{\mathbb{K}}\mathbf{L}) / \mathfrak{m}\) is a composite extension of \(\hat{\mathbf{K}}\) and \(\mathbf{L}\) over \(\mathbf{K}\) (Algebra,

Chapter VIII, § 8, Proposition 1). There exists a valuation \(w\) on \(\mathbf{L}(\mathfrak{m})\) extending \(\hat{v}\) (§ 3, no. 3, Proposition 5); the restriction \(v'\) of \(w\) to \(\mathbf{L}\) extends \(v\) . As \([\mathbf{L}(\mathfrak{m}): \hat{\mathbf{K}}]\) is finite, \(\mathbf{L}(\mathfrak{m})\) is complete with respect to \(w\) (§ 5, no. 2, Proposition 4). Now the closure of \(\mathbf{L}\) in \(\mathbf{L}(\mathfrak{m})\) is a field containing \(\hat{\mathbf{K}}\) and \(\mathbf{L}\) and hence is equal to \(\mathbf{L}(\mathfrak{m})\) . Consequently \(\mathbf{L}(\mathfrak{m})\) is identified with the completion \(\hat{\mathbf{L}}_{v'}\) and \(m\) is the kernel of the canonical mapping of \(\hat{\mathbf{K}} \otimes_{\mathbb{K}} \mathbf{L}\) onto \(\hat{\mathbf{L}}_{v'}\) . Now, by hypothesis, there exists an index \(i\) such that \(v'\) and \(v_i'\) are dependent; whence \(\mathbf{L}_{v'} = \mathbf{L}_i\) (§ 7, no. 2, Proposition 3). Thus \(n \subset m\) , which proves that \(n \subset r\) and completes the proof.

COROLLARY 1. If K is complete with respect to v and v is not improper, two valuations on L extending v are dependent.

This follows since \(\hat{\mathbf{K}}\otimes_{\mathbb{K}}\mathbf{L} = \mathbf{L}\) .

COROLLARY 2. If \(\hat{\mathbf{K}}\) or \(\mathbf{L}\) is separable over \(\mathbf{K}\) , the canonical mapping

\[
\phi : \hat {\mathbf {K}} \otimes_ {\mathbf {K}} \mathbf {L} \rightarrow \prod_ {i = 1} ^ {s} \hat {\mathbf {L}} _ {i}
\]

is an isomorphism.

The Jacobson radical of \(\hat{\mathbf{K}}\otimes_{\mathbf{K}}\mathbf{L}\) is then zero (Algebra, Chapter VIII, § 7, no. 3, Theorem 1).

Remark. Proposition 2 (b) shows that every composite extension of \(\hat{\mathbf{K}}\) and \(\mathbf{L}\) over \(\mathbf{K}\) (Algebra, Chapter VIII, § 8) is isomorphic to one of the completions \(\hat{\mathbf{L}}_t\) and that these are composite extensions no two of which are isomorphic.
% semantic-fragments: legacy-input-seam
\relax{}
